\section{Introduction}
\label{sec:intro}

Modern process operation increasingly layers data-driven components on top of the base
control system. Soft sensors infer product qualities that are too slow or too costly to
measure online; prognostic models estimate the remaining useful life of critical
equipment; and real-time optimization (RTO) continually re-tunes setpoints toward the
economic optimum. Each of these layers has, on its own, acquired rigorous uncertainty
quantification. Conformal prediction in particular supplies distribution-free,
finite-sample coverage guarantees~\citep{vovk2005,shafer2008,angelopoulos2023}, and
calibrated intervals are now practical for soft sensing, for prognostics, and for
constraint handling in optimization.

Safe operation, however, is not a property of any one layer. It is a property of the
joint outcome: the product must be in specification and the equipment must survive to its
next maintenance window. The optimizer acts on estimates, a soft-sensor interval and a
remaining-life bound, yet the event that determines safety is defined on the truths it
cannot observe. The coupling is tighter still, because the optimizer's own decision moves
every layer's operating regime at once. The reflux lever it raises for economic gain
simultaneously shifts the soft sensor away from its calibration regime and loads the pump
whose life the prognostic model is predicting. The three guarantees are therefore coupled
through the decision, and the guarantee a closed-loop operator actually needs, coverage of
the joint safety event under this coupling, does not follow from the three guarantees
taken in isolation.

A second difficulty compounds the first: conformal validity is widely held to be
incompatible with feedback. Conformal coverage rests on exchangeability between the
calibration data and the test point, and a closed loop breaks it, because the test point's
operating regime is selected by the controller using the very predictions under test. The
usual conclusion is that conformal guarantees do not survive in the loop, so conformal
methods are either confined to open-loop monitoring or deployed in feedback without a
guarantee.

This paper resolves both difficulties for the integrated system. Its contributions are:
\begin{enumerate}
\item A per-cycle \emph{composed} coverage certificate (Theorem~\ref{thm:composed}) for
the joint safety event, with terms that map one-to-one onto the three modules: the soft
sensor's miscoverage, the prognostic bound's miscoverage, a similitude-departure penalty
contributed by each module's regime shift, and a selection penalty.
\item A resolution of the feedback objection. Under the causal timing that any physically
realizable loop satisfies, in which the decision is formed before the measurement it is
conditioned against, the selection penalty is exactly zero. The conformal-selection effect
is named and bounded rather than ignored, and it is shown to be active only under a
non-causal co-selection that real loops do not exhibit.
\item A conversion of the certificate into a single convex constraint, a \emph{$\psi$-budget}
on the similitude departure, that the modifier-adaptation RTO enforces directly
(Corollary~\ref{cor:floor}). The optimizer is barred from pushing economics hard enough to
invalidate its own sensor and life guarantees, and the flagship certificate and the
operating floor reduce to the same constraint.
\item A fully reproducible closed-loop demonstration on a calibrated debutanizer twin
running the real soft-sensor, prognostic, and RTO implementations. The certified floor
$1-(\alpha_1+\alpha_2)-\varepsilon = 0.75$ is met (joint safety coverage $0.988$) when the
optimizer enforces the budget and collapses (to $0.000$) when it does not, with the
contrast driven entirely by equipment life (Figure~\ref{fig:coverage}).
\end{enumerate}

The claim is sharp: none of the three guarantees implies system safety on its own; the
composition is the certificate, and a single convex budget makes it something the optimizer
must obey. The remainder of the paper reviews the relevant
conformal, soft-sensing, prognostic, and optimization literatures and positions the
contribution; develops the framework, the certificate, and the budget; describes the
debutanizer twin and the two-arm protocol; reports the closed-loop coverage result; and
discusses conservatism, the multi-cycle horizon, and scope.
